% =============================================================================
% preamble.tex — Shared LaTeX preamble for Learning to Lens notes
% =============================================================================
% This file is \input{} by main.tex and provides all packages, macros, and
% formatting settings used across the tutorial modules.
% =============================================================================

% --- Core math and physics packages ---
\usepackage{amsmath,amssymb,amsthm}
\usepackage{mathtools}          % Extensions to amsmath (e.g., \coloneqq)
\usepackage{physics}            % Convenient notation: \dv, \pdv, \grad, \curl, etc.
\usepackage{tensor}             % Index notation: \indices{^a_b}
\usepackage{bm}                 % Bold math symbols

% --- Figures and graphics ---
\usepackage{graphicx}
\usepackage{float}              % [H] placement for figures
\usepackage{subcaption}         % Subfigures

% --- Tables ---
\usepackage{booktabs}           % Professional tables (\toprule, \midrule, \bottomrule)
\usepackage{array}

% --- Typography and layout ---
\usepackage[utf8]{inputenc}
\usepackage[T1]{fontenc}
\usepackage{microtype}          % Improved typesetting
\usepackage{enumitem}           % Customizable lists
\usepackage{fancyhdr}           % Headers and footers
\usepackage{titlesec}           % Section formatting

% --- Colors and links ---
\usepackage[dvipsnames]{xcolor}
\usepackage[
    colorlinks=true,
    linkcolor=NavyBlue,
    citecolor=ForestGreen,
    urlcolor=BrickRed,
    filecolor=Plum
]{hyperref}

% --- Code listings (for Mathematica snippets) ---
\usepackage{listings}
\lstdefinestyle{mathematica}{
    language=Mathematica,
    basicstyle=\ttfamily\small,
    keywordstyle=\color{NavyBlue}\bfseries,
    commentstyle=\color{Gray}\itshape,
    stringstyle=\color{BrickRed},
    frame=single,
    breaklines=true,
    numbers=left,
    numberstyle=\tiny\color{Gray},
    backgroundcolor=\color{gray!5},
    captionpos=b
}

% --- Bibliography ---
\usepackage[numbers,sort&compress]{natbib}

% --- Theorem environments ---
\newtheorem{theorem}{Theorem}[section]
\newtheorem{lemma}[theorem]{Lemma}
\newtheorem{proposition}[theorem]{Proposition}
\newtheorem{corollary}[theorem]{Corollary}
\theoremstyle{definition}
\newtheorem{definition}{Definition}[section]
\newtheorem{example}{Example}[section]
\newtheorem{exercise}{Exercise}[section]
\theoremstyle{remark}
\newtheorem{remark}{Remark}[section]
\newtheorem*{note}{Note}

% =============================================================================
% Custom macros — Gravitational lensing notation
% =============================================================================
% Following the conventions of Narayan & Bartelmann (1997),
% Congdon & Keeton (2018), and Saha et al. (2024).
% =============================================================================

% --- Vector quantities (bold, with arrows as needed) ---
\newcommand{\vectheta}{\boldsymbol{\theta}}     % Image position vector
\newcommand{\vecbeta}{\boldsymbol{\beta}}        % Source position vector
\newcommand{\vecalpha}{\boldsymbol{\alpha}}      % Deflection angle vector
\newcommand{\vecxi}{\boldsymbol{\xi}}            % Position in lens plane (physical)
\newcommand{\veceta}{\boldsymbol{\eta}}          % Position in source plane (physical)

% --- Angular diameter distances ---
\newcommand{\Dd}{D_{\mathrm{d}}}                 % Observer-lens distance
\newcommand{\Ds}{D_{\mathrm{s}}}                 % Observer-source distance
\newcommand{\Dds}{D_{\mathrm{ds}}}               % Lens-source distance
\newcommand{\Deff}{D_{\mathrm{eff}}}             % Effective lensing distance

% --- Critical quantities ---
\newcommand{\Sigmacr}{\Sigma_{\mathrm{cr}}}      % Critical surface mass density
\newcommand{\thetaE}{\theta_{\mathrm{E}}}        % Einstein radius (angular)
\newcommand{\RE}{R_{\mathrm{E}}}                 % Einstein radius (physical)
\newcommand{\RS}{R_{\mathrm{S}}}                 % Schwarzschild radius

% --- Lensing quantities ---
\newcommand{\kappab}{\kappa}                     % Convergence
\newcommand{\gammaL}{\gamma}                     % Shear magnitude
\newcommand{\psiL}{\psi}                         % Lensing potential
\newcommand{\alphahat}{\hat{\alpha}}             % Physical deflection angle

% --- Jacobian / magnification ---
\newcommand{\Amat}{\mathcal{A}}                  % Jacobian matrix (inverse magnification)
\newcommand{\Mmat}{\mathcal{M}}                  % Magnification tensor

% --- Cosmology ---
\newcommand{\Omegam}{\Omega_{\mathrm{m}}}        % Matter density parameter
\newcommand{\OmegaL}{\Omega_{\Lambda}}           % Dark energy density parameter
\newcommand{\Hubble}{H_0}                        % Hubble constant

% --- Physical constants ---
\newcommand{\speedoflight}{c}                     % Speed of light
\newcommand{\Gnewton}{G}                         % Newton's gravitational constant
\newcommand{\Msun}{M_{\odot}}                    % Solar mass
\newcommand{\Rsun}{R_{\odot}}                    % Solar radius

% --- Metric and tensor notation ---
\newcommand{\metric}{g_{\mu\nu}}                 % Metric tensor
\newcommand{\christoffel}[3]{\Gamma^{#1}_{\ #2 #3}} % Christoffel symbol
\newcommand{\riemann}[4]{R^{#1}_{\ #2 #3 #4}}  % Riemann tensor
\newcommand{\ricci}{R_{\mu\nu}}                  % Ricci tensor
\newcommand{\ricciS}{R}                          % Ricci scalar

% --- Commonly used shorthands ---
\newcommand{\half}{\tfrac{1}{2}}
\newcommand{\inv}[1]{{#1}^{-1}}
\newcommand{\ee}[1]{\times 10^{#1}}             % Scientific notation: 3\ee{8}
\newcommand{\order}[1]{\mathcal{O}\!\left(#1\right)} % Order of magnitude

% --- Cross-references to Mathematica files ---
% Usage: \mathlink{01a_Special_Relativity/lorentz_transforms.wl}{Lorentz transform script}
\newcommand{\mathlink}[2]{\href{../Mathematica/#1}{\texttt{#2}}}
% Usage: \solnlink{01a_Special_Relativity/problems_01a.nb}{Solutions notebook}
\newcommand{\solnlink}[2]{\href{../Solutions/#1}{\texttt{#2}}}
